% Part: first-order-logic
% Chapter: models-theories
% Section: set-theory

\documentclass[../../../include/open-logic-section]{subfiles}

\begin{document}

\olfileid[tr]{fol}{mat}{set}

\olsection{Kümeler Kuramı}

Matematiğin neredeyse tamamı kümeler kuramı içinde geliştirilebilir.
Matematiği bu kuramda geliştirmek birkaç şeyi gerektirir. Öncelikle
$\in$ bağıntısı için bir aksiyomlar kümesi gerekir. Bazen $\in$
bağıntısına birbiriyle çatışan özellikler veren çeşitli aksiyom
sistemleri geliştirilmiştir. Seçim aksiyomunu içeren Zermelo--Fraenkel
kümeler kuramı, yani $\Log{ZFC}$, bunların arasında öne çıkar: Açığa
çıktığı üzere aksiyomları, matematikçilerin ispatlamayı bekledikleri
şeylerin neredeyse tamamını ispatlamaya yeterli olduğundan, açık ara en
yaygın kullanılan ve incelenen sistemdir. Fakat bunu gösterebilmeden
önce matematikçilerin ifade etmek istedikleri şeylerin tümünü nasıl
\emph{ifade edebileceğimizi} açıklamak gerekir. Başlangıç olarak dil
hiçbir !!{constant}s ya da !!{function}s içermediğinden, belirli
kümelerden (örneğin $\emptyset$ veya $\Nat$), kümeler üzerindeki
işlemlerden (örneğin $X\cup Y$ ve $\Pow X$), hatta bağıntı ve fonksiyon
gibi küme dışındaki şeyleri içeriyor görünen yapılardan nasıl söz
edebileceğimiz ilk bakışta açık değildir.

Öncelikle ``elemanı olmak'', ilgilendiğimiz tek bağıntı değildir;
``alt kümesi olmak'' neredeyse onun kadar önemli görünür. Ancak ``alt
kümesi olmak'' bağıntısını ``elemanı olmak'' cinsinden
\emph{tanımlayabiliriz}. Bunun için kümeler kuramı dilinde, $\tuple{X,Y}$
kümeler çifti tarafından ancak ve ancak $X\subseteq Y$ olduğunda
sağlanan !!a{formula}~$!A(x,y)$ bulmalıyız. $X$, $Y$'nin alt kümesidir
ancak ve ancak $X$ kümesinin bütün !!{element}s aynı zamanda $Y$
kümesinin de !!{element}s ifadeleriyse. Dolayısıyla $\subseteq$
bağıntısını
\[
\lforall[z][(z \in x \lif z \in y)]
\]
formülüyle tanımlayabiliriz. Bundan sonra bir formülde $\subseteq$
bağıntısını kullanmak istediğimizde onun yerine bu formülü koyabiliriz
($x$ ve $y$ uygun biçimde değiştirilir; gerekiyorsa bağlı değişken~$z$
yeniden adlandırılır). Örneğin kümelerin kaplamsallığı, herhangi $x$ ve
$y$ kümeleri birbirlerini içeriyorsa aynı küme olmaları gerektiğini
söyler. Bu,
$\lforall[x][\lforall[y][((x\subseteq y\land y\subseteq x)\lif x=y)]]$
ile ya da $\subseteq$ yerine yukarıdaki tanımı koyarsak
\[
\lforall[x][\lforall[y][((\lforall[z][(z \in x \lif z \in y)] \land
    \lforall[z][(z \in y \lif z \in x)]) \lif x = y)]]
\]
ile ifade edilebilir. Bu gerçekten de $\Log{ZFC}$ aksiyomlarından biri,
\emph{kaplamsallık aksiyomudur}.

$\emptyset$ için bir !!{constant} yoktur; fakat ``$x$ boştur'' ifadesini
$\lnot\lexists[y][y\in x]$ ile belirtebiliriz. O zaman ``$\emptyset$
vardır'', $\lexists[x][\lnot\lexists[y][y\in x]]$ !!{sentence}
ifadesine dönüşür. Bu da $\Log{ZFC}$ aksiyomlarından biridir.
(Kaplamsallık aksiyomunun yalnızca bir boş küme bulunduğunu gerektirdiğine
dikkat edin.) Kümeler kuramı dilinde $\emptyset$ hakkında konuşmak
istediğimizde bunu ``boş olan bir küme vardır ve ...'' biçiminde
yazarız. Örneğin $\emptyset$ kümesinin her kümenin alt kümesi olduğunu
ifade etmek için
\[
\lexists[x][(\lnot\lexists[y][y \in x] \land \lforall[z][x \subseteq z])]
\]
yazabiliriz; elbette burada $x\subseteq z$ de kendi tanımıyla
değiştirilmelidir.

$X\cup Y$ ve $\Pow X$ gibi küme işlemlerinden söz etmek için benzer bir
yöntem kullanmalıyız. Kümeler kuramı dilinde fonksiyon simgeleri yoktur;
fakat $X\cup Y=Z$ ve $\Pow X=Y$ fonksiyonel bağıntılarını
\begin{align*}
& \lforall[u][((u \in x \lor u \in y) \liff u \in z)]\\
& \lforall[u][(u \subseteq x \liff u \in y)]
\end{align*}
ile ifade edebiliriz. Çünkü $X\cup Y$ kümesinin !!{element}s tam olarak
$X$'in ya da $Y$'nin (veya ikisinin) !!{element}s olan kümelerdir;
$\Pow X$ kümesinin !!{element}s ise tam olarak $X$'in alt kümeleridir.
Ne var ki bu, $x\cup y$ veya $\Pow x$ ifadelerini birer terimmiş gibi
kullanmamıza izin vermez; yalnızca $X\cup Y=Z$ ve $\Pow X=Y$
bağıntılarını tanımlayan bütün !!{formula}s ifadelerini kullanabiliriz.
Hatta bu bağıntıların herhangi bir zaman sağlandığını, yani birleşim ve
kuvvet kümelerinin her zaman var olduğunu henüz bilmiyoruz. Örneğin
$\lforall[x][\lexists[y][\Pow{x}=y]]$ !!{sentence}, $\Log{ZFC}$
sisteminin başka bir aksiyomudur (kuvvet kümesi aksiyomu).

Peki sıralı ikililerden veya fonksiyonlardan nasıl söz edeceğiz? Burada
sıralı ikilileri ve fonksiyonları özel türden kümeler olarak nasıl
düşünebileceğimizi açıklamalıyız. $\tuple{x,y}$ sıralı ikilisini
$\{\{x\},\{x,y\}\}$ kümesi olarak tanımlayabiliriz. Fakat daha önce
olduğu gibi, bu kümeyi adlandıran !!a{function} ekleyemeyiz; yalnızca
$\tuple{x,y}=z$, yani $\{\{x\},\{x,y\}\}=z$ bağıntısını tanımlayabiliriz:
\[
\lforall[u][(u \in z \liff (\lforall[v][(v \in u \liff v = x)] \lor
  \lforall[v][(v \in u \liff (v = x \lor v = y))]))].
\]
Bu formül, $z$ kümesinin !!{element}s olan $u$ kümelerinin tam olarak
ya yalnızca $x$'i ya da yalnızca $x$ ile $y$'yi !!{element}s olarak
içeren kümeler olduğunu söyler (başka bir deyişle bu kümeler ya
$\{x\}$ ya da $\{x,y\}$ ile özdeştir). Bunu elde ettikten sonra örneğin
$X\times Y=Z$ olduğunu da söyleyebiliriz:
\[
\lforall[z][(z \in Z \liff \lexists[x][\lexists[y][(x \in
        X \land y \in Y \land \tuple{x, y} = z)]])].
\]

$f\colon X\to Y$ fonksiyonunu $f(x)=y$ bağıntısı, yani
$\Setabs{\tuple{x,y}}{f(x)=y}$ ikilileri kümesi olarak düşünebiliriz.
O zaman bir $f$ kümesinin $X$'ten $Y$'ye bir fonksiyon olduğunu şu
koşullarla ifade edebiliriz: (a) $f\subseteq X\times Y$ bir bağıntıdır;
(b) $f$ tümeldir, yani her $x\in X$ için $\tuple{x,y}\in f$ olacak bir
$y\in Y$ vardır; (c) $f$ fonksiyoneldir, yani $\tuple{x,y},
\tuple{x,y'}\in f$ olduğunda $y=y'$ olur (çünkü fonksiyon değerleri
tektir). Dolayısıyla ``$f$, $X$'ten $Y$'ye bir fonksiyondur'' şöyle
yazılabilir:
\begin{align*}
 \lforall[u][(u \in f \lif {}] & \lexists[x][\lexists[y][(x \in X \land y \in
      Y \land \tuple{x, y} = u)]]) \land {}\\
 \lforall[x][(x \in X \lif {}] &
      (\lexists[y][(y \in Y \land \mathrm{maps}(f, x, y))] \land {}\\
& (\lforall[y][\lforall[y'][((\mathrm{maps}(f, x, y) \land
    \mathrm{maps}(f, x, y')) \lif y = y')]])).
\end{align*}
Burada $\mathrm{maps}(f,x,y)$, $\lexists[v][(v\in f\land
\tuple{x,y}=v)]$ ifadesinin kısaltmasıdır (bu !!{formula},
``$f(x)=y$'' ifadesini belirtir).

Örneğin $f\colon X\to Y$ fonksiyonunun !!{injective} olduğunu ifade
etmek de artık zor değildir:
\begin{multline*}
 f \colon X \to Y \land \lforall[x][\lforall[x'][((x \in X \land x' \in
     X \land {}] \\
 \lexists[y][(\mathrm{maps}(f, x, y) \land \mathrm{maps}(f,
        x', y))]) \lif x = x')].
\end{multline*}
Bir $f\colon X\to Y$ fonksiyonu !!{injective} ancak ve ancak $f$,
$x,x'\in X$ elemanlarını aynı $y$ elemanına eşlediğinde $x=x'$ ise.
Bu formülü $\mathrm{inj}(f,X,Y)$ ile kısaltırsak Cantor teoremi kadar
önemli bir önermeyi kümeler kuramı dilinde şimdiden ifade edebiliriz:
$\Pow X$ kümesinden $X$ kümesine !!{injective} bir fonksiyon yoktur:
\[
\lforall[X][\lforall[Y][(\Pow{X} = Y \lif
    \lnot\lexists[f][\mathrm{inj}(f, Y, X)])]].
\]

Kümeler kuramının, her tanımlayıcı özellik için bir kümenin varlığını
güvence altına alan başka bir aksiyom gerektirdiği düşünülebilir.
$!A(x)$, $x$ değişkeninin serbest olduğu bir kümeler kuramı formülü ise
\[
\lexists[y][\lforall[x][(x \in y \liff !A(x))]]
\]
!!{sentence} ifadesini ele alabiliriz. Bu !!{sentence}, !!{element}s
tam ve yalnızca $!A(x)$ özelliğini sağlayan $x$ nesneleri olan bir $y$
kümesinin bulunduğunu söyler. Bu şemaya ``küme oluşturma ilkesi''
denir. Çok yararlı görünür; ne yazık ki tutarsızdır. $!A(x)\ident
\lnot x\in x$ alırsak küme oluşturma ilkesi
\[
\lexists[y][\lforall[x][(x \in y \liff x \notin x)]]
\]
ifadesini, yani kendi kendisinin !!{element}s olmayan bütün kümelerin
kümesinin var olduğunu söyler. Böyle bir küme var olamaz; bu Russell
paradoksudur. $\Log{ZFC}$ aslında bu ilkenin kısıtlanmış ve tutarlı bir
biçimi olan ayırma ilkesini içerir:
\[
\lforall[z][\lexists[y][\lforall[x][(x \in y \liff (x \in z \land
      !A(x)))]]].
\]

\begin{prob}
\[
\lexists[y][\lforall[x][(x \in y \liff x \notin x)]] \Proves \lfalse
\]
sonucunu gösteren !!a{derivation} vererek küme oluşturma ilkesinin
tutarsız olduğunu gösterin. Önce
$(A\lif\lnot A)\land(\lnot A\lif A)\Proves\lfalse$ sonucunu göstermek
yardımcı olabilir.
\end{prob}
\end{document}
